% Generated by /paper skill — 2026-04-27
% Source: memory/paper_reference.md, .lab/results.tsv, evaluation output (1 seed, 100 maps)
% Project: PCGNN — MAP-Elites + Branching Fitness improvements

% ─── 1. COMPARISON TABLE ────────────────────────────────────────────

\begin{table}[h]
\centering
\caption{Comparison with \citet{beukman2022procedural} on 14$\times$14 maze generation
         (NoveltyNEAT baseline vs.\ our method). Results averaged over 1 seed, 100 maps.
         $\uparrow$ higher is better; $\downarrow$ lower is better.}
\label{tab:comparison}
\begin{tabular}{lccc}
\toprule
\textbf{Metric} & \textbf{Baseline~\cite{beukman2022procedural}} & \textbf{Ours} & \textbf{$\Delta$} \\
\midrule
Solvability (\%) $\uparrow$      & $100$               & $97.0$              & $-3$\,pp \\
Leniency $\uparrow$              & $0.70 \pm 0.08$     & $\mathbf{0.899}$    & $+28\%$ \\
A* Difficulty $\uparrow$         & $0.06 \pm 0.08$     & $\mathbf{0.778}$    & $+13\times$ \\
Compression Distance $\uparrow$  & $0.488 \pm 0.002$   & $0.482$             & $\approx$ \\
A* Diversity $\uparrow$          & $0.13 \pm 0.17$     & $\mathbf{0.164}$    & $+26\%$ \\
A* Edit Diversity $\uparrow$     & $0.135 \pm 0.167$   & $\mathbf{0.450}$    & $+3.3\times$ \\
\midrule
Non-L-pattern Rate (\%) $\uparrow$ & $\approx 20$\,* & $\mathbf{92}$      & $+360\%$ \\
\bottomrule
\end{tabular}
\begin{tablenotes}
  \small
  \item[*] Estimated from baseline runs; not reported in the original paper.
\end{tablenotes}
\end{table}


% ─── 2. IMPROVEMENTS (enumerate) ────────────────────────────────────

\subsection{Improvements over the Baseline}

We introduce four improvements to the original PCGNN framework~\cite{beukman2022procedural},
each targeting a specific failure mode observed in the baseline.

\begin{enumerate}

  \item \textbf{Weighted Behavior Characterization Vector (7-dim BC).}
  The original method uses visual diversity or compression distance between generators
  as the novelty distance.
  We replace this with a 7-dimensional behavior characterization vector
  $\mathbf{b} = [\rho_w,\, \ell_p,\, d_e,\, d_b,\, r,\, \delta_{A^*},\, \tau]$,
  where $\tau = n_\text{turns} / \max(1, |\pi|-2)$ encodes path tortuosity.
  Distance is computed as a weighted Euclidean norm with $w_\tau = 7.84$
  (effective $2.8\times$ amplification on the turns dimension),
  incentivising diversity in path complexity.
  This reduced the L-pattern rate from $\approx\!80\%$ to $12\%$ (non-L: 20\% $\to$ 88\%).

  \item \textbf{Mode-Collapse Penalty.}
  We observed two degenerate attractors in the baseline: diagonal-stripe walls
  and near-empty/near-full density levels.
  We add a fitness penalty
  $\mathcal{P} = \alpha\,(\bar{s}_\text{diag} + \bar{s}_\text{density})$
  with $\alpha = 0.5$, where $s_\text{diag}$ is the fraction of walls with a diagonal
  neighbour and $s_\text{density}$ penalises wall density outside $[0.15,\,0.60]$.
  Applying this penalty eliminated all diagonal-stripe maps (52 $\to$ 0)
  and reduced near-empty occurrences from 67 to 1 per 100 maps.

  \item \textbf{MAP-Elites Quality-Diversity Archive.}
  We replace the list-based novelty archive with a MAP-Elites grid
  \citep{mouret2015illuminating} of $4 \times 3 = 12$ cells
  (4 bins of path turns $\times$ 3 bins of wall density in $[0.15, 0.60]$).
  Genomes that fill an empty cell receive a $1.5\times$ fitness bonus;
  improvement earns $1.2\times$; genomes outside the density range are penalised ($0.3\times$)
  and not archived.
  Final maps are sampled only from cells with turns-bin $\geq 1$ (turns $\geq 2$),
  excluding L-pattern solutions from the output while retaining them in the archive
  to guide evolution.
  This raised the compound quality score from $0.679$ to $0.840$.

  \item \textbf{Branching Fitness Term.}
  To directly counteract the horizontal-stripe attractor induced by the tiling
  architecture, we introduce a branching fitness term
  $f_b = \operatorname{mean}\!\bigl(\nicefrac{|\mathcal{B}|}{|\mathcal{F}|}\bigr)$,
  where $\mathcal{B}$ is the set of floor cells with $\geq\!3$ walkable neighbours
  (T- and +-junctions) and $\mathcal{F}$ is all floor cells.
  This term is weighted at $w_b = 0.15$ in the composite fitness, replacing
  part of the intra-novelty weight ($0.30 \to 0.20$).
  Adding this term increased A* Difficulty from $0.06$ to $\mathbf{0.778}$ ($+13\times$)
  and A* Edit Diversity from $0.135$ to $\mathbf{0.450}$ ($+3.3\times$),
  reflecting substantially more complex and varied maze structures.

\end{enumerate}
